\chapter{METHODOLOGY}

\section{DATASET}

The dataset consists of 32,141 vehicle crossing records collected by a B-WIM
installation on a short-span bridge. Each record contains:

\begin{itemize}
    \item \textbf{Signal}: a 1,300-sample normalised strain time series covering
          one complete vehicle crossing.
    \item \textbf{Pulses}: a 1,300-sample binary array with value $1.0$ at the
          exact sample where an axle is present and $0.0$ everywhere else.
\end{itemize}

Vehicles in the dataset have between 2 and 11 axles (mean $\approx 4.5$). The
class distribution is severely imbalanced: only 0.35\,\% of all samples are
positive (axle events), giving a negative-to-positive ratio of approximately
287:1.

Records were split at the vehicle level into training (80\,\%), validation
(10\,\%), and test (10\,\%) sets using a fixed random seed, ensuring that no
vehicle appears in more than one split.

\section{BASELINE METHOD}

The baseline uses classical signal processing and serves as a performance
benchmark against which the neural network models are compared.

\begin{enumerate}
    \item \textbf{Smoothing}: A Savitzky-Golay filter \cite{SavitzkyGolay1964}
          with window length 11 and polynomial order 3 is applied to reduce
          high-frequency noise while preserving peak shape.
    \item \textbf{Peak detection}: SciPy's \texttt{find\_peaks} function
          \cite{Scipy2020} is used to locate local maxima with a minimum
          prominence of 0.05 and a minimum inter-peak distance of 5 samples.
    \item \textbf{Output}: The detected peak positions are treated as the
          predicted axle locations.
\end{enumerate}

Parameters were chosen by grid search on the validation set to maximise F1-score.

\section{CNN MODEL: 1D U-NET}

The CNN follows a 1D U-Net architecture \cite{Ronneberger2015} adapted for
sequence labelling.

\subsection{Architecture}

The encoder consists of three down-sampling stages, each comprising two
\texttt{Conv1d--BatchNorm--ReLU} blocks followed by max-pooling with stride 2.
The number of feature channels doubles at each stage (32, 64, 128). A bottleneck
layer at the lowest resolution applies two further convolutional blocks with 256
channels.

The decoder mirrors the encoder with three up-sampling stages using transposed
convolutions, each concatenated with the corresponding encoder skip connection
before two further convolutional blocks. The final layer is a $1{\times}1$
convolution that maps to a single output channel, producing a logit of shape
$[B, 1, 1300]$. No activation is applied at the output (raw logits are used
with \texttt{BCEWithLogitsLoss}).

The total number of trainable parameters is approximately 30 million.

\subsection{Input and Output}

Each signal is reshaped to $[B, 1, 1300]$ (batch, channel, time) before being
fed to the network. The target is the $[B, 1, 1300]$ pulse array. Post-training,
predicted axle positions are extracted from the output probabilities by finding
peaks with a minimum inter-peak distance of 5 samples.

\section{TCN MODEL: TEMPORAL CONVOLUTIONAL NETWORK}

\subsection{Architecture}

The TCN \cite{Bai2018} consists of nine dilated residual blocks. Block $i$ uses
dilation $d_i = 2^{i}$ (i.e., $d \in \{1, 2, 4, 8, 16, 32, 64, 128, 256\}$),
a kernel size of 3, and 64 feature channels throughout. Each block contains:

\begin{enumerate}
    \item Two \texttt{Conv1d} layers with weight normalisation, symmetric
          (non-causal) padding so that output length equals input length.
    \item \texttt{BatchNorm} and \texttt{ReLU} after each convolution.
    \item A residual connection with a $1{\times}1$ projection if channel
          dimensions differ.
\end{enumerate}

The output of the final block is passed through a $1{\times}1$ \texttt{Conv1d}
to produce a single-channel logit of shape $[B, 1, 1300]$.

\subsection{Receptive Field}

Using Equation~\ref{eq:rf} with $k=3$ and $N=9$:

\begin{equation}
    RF = 1 + 2(3-1)(2^{9}-1) = 1 + 4 \times 511 = 2{,}045 \text{ samples}.
\end{equation}

This exceeds the 1,300-sample signal length, so every output position can
potentially attend to the entire input sequence. The total number of trainable
parameters is approximately 2.7 million --- an 11-fold reduction compared to
the CNN.

\section{TRAINING PROCEDURE}

Both neural network models were trained using identical settings:

\begin{itemize}
    \item \textbf{Loss function}: \texttt{BCEWithLogitsLoss} with
          \texttt{pos\_weight}\,=\,287 (see Equation~\ref{eq:bce}), balancing the
          extreme class imbalance.
    \item \textbf{Optimiser}: AdamW \cite{Paszke2019} with initial learning rate
          $10^{-3}$, weight decay $10^{-4}$.
    \item \textbf{Learning rate schedule}: ReduceLROnPlateau (factor 0.5,
          patience 5 epochs).
    \item \textbf{Early stopping}: Training halted if validation F1 did not
          improve for 25 consecutive epochs.
    \item \textbf{Batch size}: 64.
    \item \textbf{Maximum epochs}: 100.
    \item \textbf{Hardware}: NVIDIA GPU via PyTorch 2.6 with CUDA.
\end{itemize}

The best checkpoint (highest validation F1) was saved for each model and used
for all subsequent evaluation.

\section{EVALUATION METRICS}

Evaluation is performed at the \emph{axle level}, not the sample level, to
reflect the operationally relevant quantity. Ground-truth and predicted axle
positions are first extracted from the binary/probability arrays by finding peaks
with a minimum distance of 5 samples. A predicted axle is counted as a true
positive (TP) if it falls within $\pm 5$ samples of an unmatched ground-truth
axle (greedy matching, closest first).

From these counts the following metrics are computed:

\begin{align}
    \text{Precision} &= \frac{TP}{TP + FP}, \\
    \text{Recall}    &= \frac{TP}{TP + FN}, \\
    \text{F1}        &= \frac{2 \cdot \text{Precision} \cdot \text{Recall}}
                             {\text{Precision} + \text{Recall}}.
\end{align}

Additionally, the \emph{Mean Absolute Timing Error} (MATE) measures the average
sample-level displacement of correctly matched predictions from their
ground-truth positions:

\begin{equation}
    \text{MATE} = \frac{1}{|TP|} \sum_{(p,g) \in TP} |p - g|,
\end{equation}

where $p$ and $g$ are the predicted and ground-truth axle sample indices,
respectively. A lower MATE indicates more precise timing, which translates
directly to more accurate weight estimation.